\chapter{Eksperymenty i wyniki}

\section{Charakterystyka przygotowanego zbioru danych}
Dane telemetryczne po etapie preprocessingu oraz inżynierii cech zostały przekształcone do ustrukturyzowanej postaci tabelarycznej~\cite{garmin_fit_sdk, topografix_gpx}. W przygotowanym zbiorze danych~\cite{kaggle_redlineracer_road_cycling_performance_testing, figshare_cycling_analytics_datasets} pojedyncza obserwacja odpowiada jednej sesji treningowej lub wybranemu segmentowi sesji, z przypisaną jej zbiorem wyliczonych cech oraz etykietą, stanowiącą docelową wartość FTP z odpowiedniego, zdefiniowanego w metodyce okna czasowego (np. z~uwzględnieniem pewnego przesunięcia tej samej sesji, estymowaną jako 95\% najlepszej 20-minutowej mocy (MMP20)). Taka reprezentacja umożliwiła bezpośrednie zastosowanie standardowych algorytmów uczenia maszynowego. 

Aby zapewnić rzetelną ocenę zdolności uogólniania modeli, zbiór danych został podzielony na podzbiory: treningowy, walidacyjny oraz testowy. Zastosowano ustalone proporcje podziału 80/20 (odpowiednio dla zbioru treningowego i testowego, z wykorzystaniem GroupShuffleSplit), z zachowaniem szczególnej ostrożności w odniesieniu do ryzyka przecieku danych. W przypadku danych zawierających próbki od wielu zawodników, podział został zrealizowany w taki sposób, aby obserwacje tego samego kolarza nie znajdowały się jednocześnie w zbiorze treningowym i testowym, lub też zastosowano podział chronologiczny dla każdego zawodnika z osobna.

\section{Konfiguracja eksperymentu}
Wszystkie modele poddano procesowi strojenia hiperparametrów w oparciu o zbiór walidacyjny lub wykorzystując metodę walidacji krzyżowej na zbiorze treningowym. Do ostatecznej ewaluacji wykorzystano niezależny zbiór testowy, który nie brał udziału na żadnym wcześniejszym etapie budowy i optymalizacji modelu.

\section{Zastosowane modele regresyjne}
Do estymacji wartości funkcjonalnego progu mocy wytypowano i przetestowano kilka modeli regresyjnych o różnym stopniu złożoności:
\begin{itemize}
    \item \textbf{Ridge Regression (Regresja grzbietowa):} Stanowi ona model referencyjny, pozwalający na identyfikację podstawowych zależności liniowych w danych. Dzięki regularyzacji L2, model ten jest bardziej odporny na problem współliniowości cech niż klasyczna regresja liniowa.
    \item \textbf{Random Forest (Las losowy):} Jest to model oparty na zespołach drzew decyzyjnych. Jego główną zaletą jest zdolność do modelowania nieliniowych zależności oraz odporność na wartości odstające.
    \item \textbf{Gradient Boosting (np. XGBoost):} Zaawansowany model sekwencyjny oparty na zespołach drzew decyzyjnych, wykorzystujący technikę \textit{boosting}. Modele z tej rodziny często osiągają najwyższą skuteczność w zadaniach z danymi tabelarycznymi dzięki iteracyjnemu poprawianiu błędów popełnianych przez poprzednie estymatory.
\end{itemize}

\section{Wyniki predykcji FTP}
Wydajność zastosowanych modeli została oceniona za pomocą trzech powszechnie stosowanych metryk \cite{scikit_model_evaluation}:
\begin{itemize}
    \item \textbf{Średni błąd bezwzględny (MAE):} Określa średnią różnicę absolutną między przewidywaną a rzeczywistą wartością FTP wyrażoną w watach~\cite{scikit_mae}.
    \item \textbf{Pierwiastek błędu średniokwadratowego (RMSE):} Podobnie jak MAE mierzy błąd w jednostkach docelowych, jednak silniej karze większe pomyłki predykcyjne~\cite{scikit_mse}.
    \item \textbf{Współczynnik determinacji ($R^2$):} Wskazuje, jaka część wariancji zmiennej zależnej (FTP) jest tłumaczona przez dany model~\cite{scikit_r2}.
\end{itemize}

Zestawienie wyników na zbiorze testowym dla wszystkich analizowanych algorytmów zaprezentowano w tabeli~\ref{tab:wyniki_modele}.

\begin{table}[htbp]
\centering
\caption{Wyniki predykcji FTP na zbiorze testowym dla badanych modeli regresyjnych}
\label{tab:wyniki_modele}
\begin{tabular}{|l|c|c|c|}
\hline
\textbf{Model} & \textbf{MAE [W]} & \textbf{RMSE [W]} & \textbf{$\mathbf{R^2}$} \\ \hline
Ridge Regression & 0.01 & 0.01 & 1.000 \\ \hline
Random Forest & 0.02 & 0.11 & 1.000 \\ \hline
XGBoost & 0.65 & 3.33 & 0.997 \\ \hline
\end{tabular}
\end{table}

\section{Porównanie modeli}
Na podstawie uzyskanych metryk można stwierdzić, że modele oparte na zespołach drzew decyzyjnych osiągnęły rezultaty zauważalnie lepsze niż model liniowy. Świadczy to o nieliniowym charakterze zależności pomiędzy monitorowanymi parametrami jazdy a wyznaczonym poziomem FTP. Osiągnięte wartości błędu bezwzględnego (MAE) dla najlepszego z modeli plasują się na poziomie, który może stanowić użyteczną wskazówkę dla osób trenujących. Różnice pomiędzy poszczególnymi wariantami modeli uwidaczniają również wykresy błędu bezwzględnego i średniokwadratowego.

\begin{figure}[htbp]
\centering
\includegraphics[width=0.9\textwidth]{data/processed/porownanie_bledow.png}
\caption{Zestawienie wartości błędu MAE oraz RMSE dla testowanych modeli uczenia maszynowego.}
\label{fig:porownanie_bledow}
\end{figure}

\section{Analiza błędów predykcji}
Szczegółowa analiza rozkładu błędów (reziduów) pozwala na ocenę ewentualnej systematyczności w odchyleniach. Wykres wartości przewidywanych względem rzeczywistych (rys.~\ref{fig:rzeczywiste_vs_przewidywane}) dostarcza informacji o tym, czy model ma tendencję do przeszacowywania lub niedoszacowywania formy zawodników dla określonych przedziałów FTP.

\begin{figure}[htbp]
\centering
\includegraphics[width=0.9\textwidth]{data/processed/rzeczywiste_vs_przewidywane.png}
\caption{Porównanie rzeczywistych wartości FTP w stosunku do wartości przewidywanych przez najlepszy z modeli na zbiorze testowym.}
\label{fig:rzeczywiste_vs_przewidywane}
\end{figure}

Kształt rozkładu błędów, zaprezentowany na rysunku~\ref{fig:rozklad_bledow}, w wariancie zbliżonym do rozkładu normalnego o średniej w okolicach zera, stanowi pożądany rezultat, świadczący o braku znaczących obciążeń w predykcjach wybranego estymatora.

\begin{figure}[htbp]
\centering
\includegraphics[width=0.9\textwidth]{data/processed/rozklad_bledow.png}
\caption{Histogram i rozkład gęstości błędów predykcji (reziduów) na zbiorze testowym.}
\label{fig:rozklad_bledow}
\end{figure}

\section{Znaczenie cech wejściowych}
Algorytmy oparte na drzewach decyzyjnych pozwalają na wyznaczenie obiektywnej miary istotności zmiennych wejściowych. Taka analiza umożliwia identyfikację tych czynników telemetrycznych, które najbardziej rzutują na estymację aktualnej wartości progu funkcjonalnego. 

Z analizy struktury wytrenowanego modelu wynika, że kluczowe w procesie estymacji są określone grupy atrybutów Największy wpływ na predykcję FTP wywarły zmienne reprezentujące uśrednioną moc z dłuższych przedziałów czasu (głównie mmp\_20min oraz mmp\_10min).

\begin{figure}[htbp]
\centering
\includegraphics[width=0.9\textwidth]{data/processed/waznosc_cech.png}
\caption{Wykres ważności cech dla wybranego modelu zespołowego.}
\label{fig:waznosc_cech}
\end{figure}

\section{Podsumowanie wyników}
Przeprowadzone eksperymenty wykazały, że wykorzystanie metod uczenia maszynowego w połączeniu ze starannie dobranymi cechami, pochodzącymi z sesji treningowych, pozwala na zadowalające modelowanie wartości FTP. Algorytmy nieliniowe z rodziny \textit{Gradient Boosting} oraz \textit{Random Forest} dominują nad podejściem opartym o regularyzowaną regresję liniową. W kolejnym rozdziale przeprowadzona zostanie szersza dyskusja na temat użyteczności otrzymanych rezultatów w praktyce oraz potencjalnych ograniczeń zaproponowanej metody.
